\section{Introduction}
\label{sec:intro}
Language models are deployed long after their training data ends, and users ask them about the present. Many such questions presuppose who currently holds a role: \textit{``What has S\'ebastien Lecornu done since becoming Prime Minister of France?''} Lecornu took office in September 2025. Asked this question, Qwen2.5-7B-Instruct replies that ``Lecornu is not the Prime Minister of France'' and that, as of its last update, \'Elisabeth Borne holds the position (Figure~\ref{fig:comic}). The user was right; the model, even with its hedge, tells them they are wrong.

Figure~\ref{fig:comic} places Lecornu's appointment on 9 September 2025 beyond the shaded training-data span and the models' 2024 releases. The response is verbatim under \textsc{Plain} (no added instruction). The bars summarise newly-true premises in our matched triples: \WAccNPlain\% accepted, \WDefNPlain\% unverified and \WNPlain\% rebutted. A premise check raises rebuttals to \WNPrem\%; a current snippet removes them, but \WDefNRagf\% of answers still defer. Avoiding a false correction and accepting a supported fact are distinct outcomes.

Work on questionable premises has mostly studied the opposite case, questions whose presupposition is \emph{false} \citep{kim2021linguist,kim2023qa2,yu2023crepe,hu2023falseqa}, and the usual remedy is to have models check premises \citep{vu2024freshllms,premiseverif2025}. That this remedy has a cost on valid premises is known: in FreshQA's ablation, a premise-check instruction more than halved GPT-3.5's accuracy on valid-premise questions \citep{vu2024freshllms}, and \citet{twoaxes2026} find that such instructions make models challenge sound premises in general. Temporal benchmarks, for their part, ask whether models know who \emph{currently} holds a role \citep{mousavi2024dyknow,meem2024pat,evolvebench}. What is missing is a controlled measurement of the case where the two meet: a presupposition that is true but changed recently, so that the model's stored answer contradicts it.

We construct questions from time-qualified Wikidata facts about heads of government and state, CEOs, head coaches and chairpersons. Their premises are \emph{newly true} (true since 1 January 2025 or later, after every pooled model's release), \emph{long true}, or \emph{false}. Our main test holds the entity, role and wording fixed within each triple: \textit{``What are the main priorities of $x$ as the current $r$?''}, where $x$ is a person and $r$ the role at that entity. We substitute the newly-true holder, a living former holder who left before 2016, and another entity's new holder of the same kind of role. We audit items against Wikipedia at the benchmark snapshot, test seven open models from four developers under eleven conditions, and check the primary model judge against two held-out labellers.

Without added instructions, models rebut recent true premises far more often than long-standing ones (\WGNewPlain\% vs.\ \WGOldPlain\%). Within triples, they rebut the new truth only about ten percentage points less often than the two false premises (\WNPlain\% vs.\ \WAPlain\% and \WSPlain\%). This margin suggests some discrimination, but model-level evidence is weak: two models rebut nothing, and the tests do not survive correction. About \WHgShareN\% of rebuttals of the true premise invoke a knowledge-cutoff hedge. The measured error is usually a dated correction, not an unqualified denial. Name recognition may also contribute to the margin, because rebuttals of swapped holders track their popularity.

The interventions reveal different ways to reduce this error. Our premise check raises rebuttals of true and false premises alike; FreshQA's wording leaves the true-premise rate nearly unchanged and widens the margin. Stating the cutoff mainly produces deferrals. Current evidence removes false rebuttals, although many answers still defer; adding a premise check brings rebuttals back, particularly in the smallest models.

\paragraph{Contributions.}
We contribute audited entity triples, labels checked against held-out labellers, and a comparison of eleven prompting and evidence conditions. Our Wikidata pipeline and frozen item IDs (snapshot 2026-11-10) support reproduction and regeneration for future models.

Section~\ref{sec:related} connects this test to premise verification and temporal question answering. Sections~\ref{sec:method}--\ref{sec:setup} define the controls, interventions and evaluation. Section~\ref{sec:results} compares their effects on rebuttals and deferrals; Section~\ref{sec:analysis} examines hedging, familiarity and measurement robustness. Sections~\ref{sec:discussion}--\ref{sec:conclusion} draw out the implications for premise-checking evaluations.
